\documentclass[11pt,a4paper]{article}
\usepackage[T1]{fontenc}
\usepackage[utf8]{inputenc}
\usepackage{lmodern}
\usepackage[a4paper,margin=25mm]{geometry}
\usepackage{amsmath,amssymb,amsthm,mathtools}
\usepackage{booktabs,longtable,array,tabularx}
\usepackage{microtype}
\usepackage{xcolor}
\usepackage{enumitem}
\usepackage{fancyhdr}
\usepackage{xurl}
\usepackage[unicode,breaklinks,colorlinks=true,linkcolor=blue!45!black,citecolor=blue!45!black,urlcolor=blue!45!black]{hyperref}
\hypersetup{pdftitle={Sign recovery and centered norms for explicit ZFC consistency equations},pdfauthor={AI-assisted research contribution prepared for ProveIt}}
\setlength{\emergencystretch}{2em}
\setlength{\parskip}{0.35em}
\setlength{\parindent}{1.2em}
\setlist{topsep=4pt,itemsep=2pt,parsep=1pt}
\allowdisplaybreaks[2]
\numberwithin{equation}{section}
\newcommand{\ZFC}{\mathrm{ZFC}}
\newcommand{\Con}{\operatorname{Con}}
\newcommand{\Z}{\mathbb Z}
\newcommand{\N}{\mathbb N}
\newcommand{\Zpos}{\mathbb Z_{>0}}
\newcommand{\snapshot}{\texttt{beb0192cb8dc36f78c369c080f2adfbc1887a486}}
\newtheorem{theorem}{Theorem}[section]
\newtheorem{lemma}[theorem]{Lemma}
\newtheorem{proposition}[theorem]{Proposition}
\pagestyle{fancy}
\fancyhf{}
\fancyhead[L]{\small Explicit consistency equations}
\fancyhead[R]{\small Sign recovery and centered norms}
\fancyfoot[C]{\thepage}
\setlength{\headheight}{14pt}
\title{Sign recovery and centered norms\\for explicit ZFC consistency equations}
\author{Research contribution prepared for the ProveIt project}
\date{23 October 2026}

\begin{document}
\maketitle
\begin{abstract}
We give explicit improvements to several Diophantine equations calibrated
to the syntactic consistency of ZFC. The main short equation has thirteen
strictly positive integer unknowns, exact total degree $114$, and an
$83$-operation division-free evaluator. Centering an algebraic norm at a
known large part of a Pell root replaces two root witnesses by one
centered witness with lower degree. A separate common-sign argument gives a quartic in forty-two
unrestricted integers: two independent positive representations can be
replaced by one positive representation of their product, because a larger
existing square scale lets the Pell indices recover the missing sign.
A proved exchange of scale guards followed by exact eliminations gives a
sextic in thirty-five integers and an octic in thirty integers. A companion sparse packing optimization gives nine
positive unknowns, degree $1\,834\,516$, and $453$ operations.

Each equation comes with a literal polynomial or arithmetic circuit,
fixed coefficient payloads, a two-way ordinary mathematical reduction,
and finite verification programs. The comparisons use the pinned ProveIt
revision stated below. The native computation compiler and K215 calibration
are inherited. We do not claim a globally simplest equation, an
independent literature priority result, or a new end-to-end formal proof.
These are concrete local improvements, not an established major
breakthrough in Diophantine representation.
\end{abstract}

\clearpage
\tableofcontents
\clearpage

\section{Results, comparisons, and the exact logical claim}
\label{sec:overview}

The most readable new short equation is $P_{13}=0$ in
Section~\ref{sec:positive13}. It uses thirteen positive inputs and four
\emph{fixed} integer coefficients. Its degree is $114$ and its evaluator
uses $83$ binary additions, subtractions, and multiplications, including
all scalar multiplications and squarings. The lowest-degree new equation
is the quartic $Q_{42}=0$ in Section~\ref{sec:quartic42}; all forty-two of
its unknowns range freely over $\Z$.

The comparison source is the repository
\texttt{VladimirReshetnikov/ProveIt}, directory
\path{Computability/DiophantineConsistency/con-zfc-diophantine-forms},
at the commit
\begin{center}\snapshot.\end{center}
This includes the six wave-eight packages numbered 54 through 59. In
particular, the older collection index alone is not the comparison
frontier. The relevant late package READMEs and exact predecessor objects
are retained or pinned in the accompanying archive
\cite{repo57,repo58,repo59}.

\begin{table}[htbp]
\centering
\small
\begin{tabular}{@{}llrrl@{}}
\toprule
Domain & Construction & Unknowns & Degree & Binary operations\\
\midrule
$\Zpos$ & Source 59 short adjunct & 13 & 126 & 91\\
$\Zpos$ & Source 57 short form & 13 & 132 & 83\\
$\Zpos$ & \textbf{Centered norm, this work} & \textbf{13} & \textbf{114} & \textbf{83}\\
\addlinespace
$\Z$ & Source 57 quartic & 45 & 4 & Not optimized\\
$\Z$ & \textbf{Common-sign quartic} & \textbf{42} & \textbf{4} & Not optimized\\
$\Z$ & Archived sextic, sources 27/29 & 42 & 6 & Not optimized\\
$\Z$ & \textbf{New sextic companion} & \textbf{35} & \textbf{6} & Not optimized\\
$\Z$ & Archived octic, sources 27/29 & 36 & 8 & Not optimized\\
$\Z$ & \textbf{New octic companion} & \textbf{30} & \textbf{8} & Not optimized\\
\addlinespace
$\Zpos$ & Source 59 sparse form & 9 & $1\,910\,404$ & 473\\
$\Zpos$ & \textbf{Repacked companion} & \textbf{9} & $\boldsymbol{1\,834\,516}$ & \textbf{453}\\
\bottomrule
\end{tabular}
\caption{Exact comparisons within each domain and degree or variable-count
regime. The old sextic and octic are historical low-degree baselines, not
the fewest-variable unrestricted-integer equations in the collection.}
\label{tab:results}
\end{table}

There are additional useful tradeoffs in the repository. For example,
source 58 has a fifteen-positive-variable equation of degree $24$, and
source 57 has twelve-positive-variable equations of degree $242$ and
$258$ \cite{repo57,repo58}. This paper does not improve those points.
The different domains in Table~\ref{tab:results} must not be conflated:
thirteen positive variables and thirteen unrestricted integer variables
are different representation problems. No nine-unrestricted-variable
claim is made here.

\subsection{The theorem being proved}

Fix the repository's effective presentation of ZFC and write
$\operatorname{Ref}_{\ZFC}(p)$ for its primitive recursive proof-checking
predicate for a refutation. Define the syntactic consistency statement by
\[
 \Con(\ZFC)\ \equiv\ \neg\exists p\in\N\,
                                \operatorname{Ref}_{\ZFC}(p).
\]
For every polynomial $P$ in Table~\ref{tab:results} introduced by this
work, and its stated domain $D$, the claimed mathematical theorem is
\begin{equation}
 \ZFC\ \vdash\
 \left(\exists\boldsymbol{x}\in D^n\;P(\boldsymbol{x})=0\right)
 \longleftrightarrow
 \left(\exists p\in\N\;\operatorname{Ref}_{\ZFC}(p)\right).
 \label{eq:logical-claim}
\end{equation}
Negation gives precisely the requested equivalence
\begin{equation}
 \ZFC\ \vdash\
 \left(\neg\exists\boldsymbol{x}\in D^n\;P(\boldsymbol{x})=0\right)
 \longleftrightarrow\Con(\ZFC).
 \label{eq:consistency-claim}
\end{equation}
Neither implication assumes that ZFC is consistent. A proof of
\eqref{eq:consistency-claim} is also not a proof in ZFC that $P$ has no
zero. The latter would prove the consistency statement itself.

The new steps are arithmetic: integer polynomial identities, Pell
recurrences and classification, divisibility, elementary bounds, and
finite coefficient constructions. They are ordinary mathematics
formalizable in ZFC. Their composition uses the expressly identified
native compiler and literal K215-to-refutation calibration. Thus the new
constructions remain based on the inherited K215 endpoint. The quartic,
sextic, octic, and thirteen-variable equation do not need a Jones or Sun
variable-reduction theorem. The nine-variable companion retains its
Jones ratio interface.

\subsection{What ``simpler'' measures here}

We report the domain, number of quantified unknowns, exact total degree,
arithmetic operation count when optimized, and fixed coefficient size
separately. An arithmetic register is an abbreviation, not an extra
existential unknown. Every exponent in an emitted polynomial circuit is
a fixed numeral, and every power is expanded to charged multiplication
gates. Constants are leaves; reusing a computed register is free. This
counts arithmetic operations, not bit operations on the very large
coefficients or witnesses.

The thirteen-variable result strictly improves both displayed predecessor
points in degree without increasing either operation count. The
common-sign construction reduces the number of integer unknowns at fixed
quartic degree. The nine-variable result improves its predecessor by
$75\,888$ degrees and $20$ operations. None of these finite comparisons
establishes a global minimum. MRDP establishes existence of appropriate
Diophantine representations \cite{matiyasevich,jones}; it does not certify
an optimal explicit instance. Also, current unrestricted-integer
variable-count results must not be summarized by the older eleven-variable
record: Sun's September 2026 preprint states a ten-variable undecidability
theorem \cite{sun2026}. That universal theorem is background, not an
additional reduction silently applied to the equations here.

\section{Fixed coefficients and inherited dependencies}
\label{sec:binding}

A display containing four unexplained coefficient symbols would specify
only a family. Here those symbols are bound to complete finite byte
strings, together with deterministic reconstruction programs and their
input data. The fixed integers are unsigned, minimally encoded in
little-endian order. The compressed payloads carry no mathematical
ambiguity: their decompressed sizes and cryptographic digests are checked.
The native transition table and tile definitions used to reconstruct them
are included.

\begin{table}[htbp]
\centering\small
\begin{tabular}{@{}lrrrr@{}}
\toprule
Whole native tuple & $\beta$ bits & $K_0$ bits & $M_C$ bits & $M_F$ bits\\
\midrule
Compact 44: $P_{13}$ and nine & 590069 & 592884 & 590069 & 590070\\
Source 34: $Q_{42}$ and companions & 1064921 & 1065043 & 1064921 & 1064922\\
\bottomrule
\end{tabular}
\caption{Bit lengths of the fixed, coherent native tuples. The full
manifests and all four numeral payloads accompany each construction.}
\label{tab:ports}
\end{table}

For the compact tuple, $\beta=2^{590069}-1$ and the four raw integer
encodings occupy $295\,388$ bytes. For the quartic tuple,
$\beta=2^{1064921}-1$ and they occupy $532\,479$ bytes. The fourth port
$M_F$ is the \emph{already shifted} source mask in both presentations.
The exponent of $\beta$ alone is not an instance specification, and
individual ports from different tuples must not be mixed. The
nine-variable final evaluator has twelve fixed leaves derived from its
native tuple; their encodings total $2\,217\,161$ raw bytes and $47\,364$
gzip bytes. Those final leaves are supplied too.

The dependency chain has three mathematical layers. At the bottom,
a literal K215 halting computation is equivalent to a refutation in the
fixed ZFC calculus. The native window compiler and its cofinal padding
turn that halting computation into positive solutions of the bound
arithmetic source system and decode every source solution back to such
a computation. Finally, the new polynomial transformations in this paper
preserve existence of solutions in both directions. The archive retains
the ordinary proofs used at each interface and records exact repository
paths, commit pins, Git blob identities and SHA256 digests. The detailed
reading map is
\path{quartic42/dependencies/README.md}.

The included formal endpoint is
\begin{center}
\small\texttt{KEndpoint215.blank\_halts\_iff\_zfc\_refutation}.
\end{center}
It identifies blank halting with the repository's
\texttt{BProv ZFCax\_s [] Form.fBot}. Its complete transitive Lean/Rocq
development is not rebuilt in this archive. More importantly, an inherited
machine theorem is not a new kernel certificate for the literal equations
introduced here. We provide ordinary mathematical reductions and exact
finite data checks; we do not provide an exported object-language ZFC
derivation or a fresh end-to-end proof-assistant build. This boundary is
part of the result, not an unstated hypothesis that ZFC is consistent.

\section{A root-centered equation in thirteen positive integers}
\label{sec:positive13}

The construction in this section has thirteen strictly positive integer
unknowns, exact total degree $114$, and $83$ live binary arithmetic
operations: $46$ multiplications, $19$ additions, and $18$ subtractions.
Squaring and multiplication by a fixed scalar are charged. Intermediate
registers are polynomial abbreviations and do not add unknowns.

The new step merges the first Pell root and the auxiliary Pell ordinate
using a centered two-root norm. Keeping the ordinary strong coordinate
allows direct recovery of both deleted roots. This gives a degree
improvement over the pinned thirteen-positive constructions of degree
$126$ in $91$ operations and degree $132$ in $83$ operations. These are
comparisons with the named repository constructions, without a claim of
global minimality or literature priority.

\subsection{The fixed equation}

Use the entire compact44 native tuple, with
$\beta=2^{590069}-1$. The four payloads and their exact hashes are fixed by
\texttt{COEFFICIENT\_BINDING.json}; the aliases are
$\beta\mapsto\beta$, $K_0\mapsto K$, $M_C\mapsto M_C$, and
$M_F\mapsto M_{F,\mathrm{source}}$.
The last coefficient is already shifted:
$M_{F,\mathrm{source}}=M_{F,\mathrm{native}}+2^{590069}-1$.
The whole tuple has $295388$ raw bytes. Its included proof supplies the
same native blank-machine compiler interface, with the ordinary strong
witness obtained from a canonical normalized one by
$i_{\mathrm{ordinary}}=\Delta i_{\mathrm{normalized}}$.

Supply precisely
\[
 J,F,\alpha,z,f,i,T,s,w,\eta,\zeta,\gamma,t
 \quad\in\mathbb Z_{>0}.
\]
Define the following polynomial abbreviations:
\begin{align*}
 q&=\beta J+1,& C&=q-F-\alpha,\\
 X&=wq,&Y&=sq^3,&E&=XY,& k&=\eta+\zeta,\\
 c&=kY+\eta,&a&=E+Y,&H&=4a+3,&\Delta&=a^2+H,\\
 \ell&=EkY,& A_1&=1+\ell(\ell+k),&G&=1+\Delta c^2,&b_0&=X+ac.
\end{align*}
The packing and transport expressions are
\begin{align*}
 R&=\bigl((q-1)(q-F)+\alpha\bigr)(q^2-1)
       +(M_C+qM_F)J,\\
 N_t&=(K_0+w)C+q-F-z(q-1).
\end{align*}
For the absorbed main root put
\begin{align*}
 U&=4q^4,& r&=k-R-1,\\
 D_*&=Ub_0+H(U\gamma-r),&
 N_{\mathrm{main}}&=D_*^2-U^2G.
\end{align*}
The strong and auxiliary quantities are
\begin{align*}
 I&=ic^2,&Q&=\Delta(f^2-1),&d&=Q-1,&N_s&=I^2-d,\\
 V&=c(Tf-1)-Rf^2,&E_{\mathrm{aux}}&=QV^2-1,\\
 x&=t+\ell,&\delta&=x^2-A_1,&
 N&=(d\delta-E_{\mathrm{aux}})^2-4dA_1E_{\mathrm{aux}}.
\end{align*}
The proposed equation is
\begin{equation}
 \boxed{\displaystyle
 P_{13}=N_t(1+N_{\mathrm{main}})N_s(1+N)-1=0.}
 \label{eq:p13-main}
\end{equation}
Every displayed power has a fixed nonnegative integer exponent.
The included division-free JSON circuit expands every abbreviation.
The coordinate $z$ here is the old positive \emph{product} transport
coordinate. It equals $z_{\mathrm{shift}}+1$ for the coordinate used in a
shifted equation $(K_0+w)C=F+z_{\mathrm{shift}}(q-1)$.

With the fixed tuple, ordinary mathematics formalizable in ZFC proves
\[
 \exists\boldsymbol v\in\mathbb Z_{>0}^{13}\;
 P_{13}(\boldsymbol v)=0
 \quad\Longleftrightarrow\quad
 \text{there is a finite syntactic ZFC refutation}.
\]
Negating this equivalence gives the required equivalence between the
absence of positive zeros and $\operatorname{Con}(\mathrm{ZFC})$.
The native compiler and the literal K215 calibration are inherited
theorems; no new end-to-end proof-assistant certificate is claimed.

\subsection{The centered two-root identity}

For arbitrary indeterminates $d,x,A,e$ there are exact polynomial identities
\begin{align}
 (d(x^2-A)-e)^2-4dAe
 &=\bigl(d(x^2+A)-e\bigr)^2-4d^2x^2A \notag\\
 &=d^2(x^2-A)^2-2d(x^2+A)e+e^2.
 \label{eq:p13-norm-identity}
\end{align}
The first expression saves one multiplication in the evaluator. If
$d,x,A$ are positive integers and it vanishes, the middle expression gives
the rational number
\[
 \tau_0=\frac{d(x^2+A)-e}{2dx},\qquad \tau_0^2=A.
\]
A rational square root of an integer is an integer. Therefore
$\tau_0\in\mathbb Z$ and
\begin{equation}
 d(x-\tau_0)^2=e.
 \label{eq:p13-second-root}
\end{equation}
This direct argument allows coincident square classes. It requires no
unproved independence of formal radical signs and no auxiliary rational
Pell saturation lemma.

The particular shift $x=t+\ell$ gives
\begin{equation}
 \delta=t^2+\ell(2t-k)-1.
 \label{eq:p13-centering}
\end{equation}
This cancellation lowers its degree from the apparent $22$ to $12$.
The converse will prove positivity of the new witness
$t=\tau-\ell+y$; positivity is not assumed to be preserved by an arbitrary
affine change of an old tuple.

\subsection{The unit guard and the two restored roots}

At a positive zero of \eqref{eq:p13-main}, each of the four integer factors
is $1$ or $-1$. Since $4\mid U$, the residue of $N_{\mathrm{main}}$ modulo
$4$ is the square $D_*^2$. The residue of $N$ modulo $4$ is also a square.
Neither residual can be $-2$. Thus
$N_{\mathrm{main}}=N=0$ and their two factors are $+1$.

For every integer $a$, $\Delta=a^2+4a+3$ is $0$ or $3$ modulo $4$.
Hence $Q=\Delta(f^2-1)$ is $0$ or $1$ modulo $4$.
It follows that $N_s=1+I^2-Q$ is $0$, $1$, or $2$ modulo $4$ and cannot
be $-1$. Consequently
\begin{equation}
 Q=I^2,\qquad N_t=1,\qquad N_{\mathrm{main}}=N=0.
 \label{eq:p13-guard-conclusions}
\end{equation}
Here $c>1$ and $I=ic^2>1$, so $Q>1$ and $d>0$. In particular the
apparently troublesome case $f=1$ is excluded by the equations themselves.
Also $x=t+\ell>0$ and $A_1>0$.

If $V=0$, then $E_{\mathrm{aux}}=-1$, and
\[
 N=(d\delta+1)^2+4dA_1>0,
\]
a contradiction. Therefore $V\ne0$ and $E_{\mathrm{aux}}>0$.
Applying \eqref{eq:p13-norm-identity} and
\eqref{eq:p13-second-root} gives integers
$\tau_0$ and $y_0=x-\tau_0$ with
\[
 \tau_0^2=A_1,\qquad dy_0^2=E_{\mathrm{aux}}.
\]
Thus the positive integers $\tau=|\tau_0|$ and $y=|y_0|$ satisfy
\begin{equation}
 \tau^2-\ell(\ell+k)=1,\qquad
 QV^2-(Q-1)y^2=1,\qquad
 (ic^2)^2=\Delta(f^2-1).
 \label{eq:p13-restored-roots}
\end{equation}
The identity
$\ell(\ell+k)=(XY^2)(XY^2+1)k^2$ identifies the first original Pell norm.
Together with \eqref{eq:p13-guard-conclusions}, these are precisely the
ordinary-strong fourteen-coordinate absorber equations, with all retained
coordinates still positive. Neither an index quotient nor its sign has
yet been assumed.

\subsection{Early index recovery and restoration of the compiler}

The original absorber proof is supplied verbatim as
\texttt{upstream/FOURTEEN\_ABSORPTION\_PROOF.md}; its full source identity
is recorded in the manifest. We restate its proof order because the early
index bound is necessary for the ordinary strong-rank lemma.

The equality $N_t=1$ implies $C>0$: if $C\le0$, positivity of $K_0+w$ gives
\[
 N_t\le q-F-z(q-1)=1-F-(z-1)(q-1)\le0.
\]
The authentic mask bounds then give
\begin{equation}
 q\ge16,\quad 3q<R<q^4-2,\quad E\ge q^4>R+2,\quad
 kY<c<k(Y+1),\quad a+2<2E.
 \label{eq:p13-outer-bounds}
\end{equation}
From $N_{\mathrm{main}}=0$, $g=D_*/U$ is a rational square root of the
integer $G$, hence an integer. It is positive. Indeed, $H<5a$, $b_0>akY$,
$r<k$, and $\gamma>0$ imply
\[
 D_*>Ub_0-Hk>ak(UY-5)>0.
\]
Thus $g^2-\Delta c^2=1$. Reduction of $D_*=Ug$ modulo $U$ gives
\begin{equation}
 U\mid H(k-R-1).
 \label{eq:p13-absorbed-congruence}
\end{equation}

Put $A_0=a+2$ and $P_0=2XY^2+1$. The usual positive Pell sequences give
indices $p,n\ge1$ such that
\[
 g=\chi_{A_0}(p),\quad c=\psi_{A_0}(p),\qquad
 \tau=\chi_{P_0}(n),\quad k=2\psi_{P_0}(n).
\]
Since $P_0>A_0$ and $c>k$, monotonicity implies $p>n$.
The first recurrence gives $k\equiv2n\pmod{4XY^2}$, and
$U\mid4XY^2=4ws^2q^7$. Moreover $H$ is odd and $H\equiv3\pmod q$;
if $3\mid q$, then $H\equiv3\pmod{27}$. Hence
$\gcd(U,H)\le3$. Defining $\mu=U/\gcd(U,H)$, we obtain
\[
 \mu\ge\frac{4q^4}{3}>R+1,\qquad
 2n\equiv R+1\pmod\mu,
 \qquad n\ge\frac{R+1}{2}\ge25.
\]
The lower bound precedes recovery of $p=R$. It gives $p\ge26$, and Pell
growth therefore proves
\begin{equation}
 c>A_0\Delta^2,\qquad c>2p,\qquad c>2R.
 \label{eq:p13-rank-hypotheses}
\end{equation}
For example,
$c\ge(2A_0-1)^{p-1}\ge A_0^{25}>A_0(A_0^2-1)^2$;
also $c\ge\psi_{A_0}(2)=2A_0>2R$.

The ordinary strong-rank lemma in the supplied
\texttt{ORDINARY\_BLANK\_PROOF.md}, Section~4, now applies. It yields
\begin{equation}
 \begin{gathered}
 f=\chi_{A_0}(m),\quad p\mid m,\quad c\mid m,\quad m\ge c>2p,\\
 I=\Delta\psi_{A_0}(m),\qquad c^2\mid(f^2-1).
 \end{gathered}
 \label{eq:p13-strong-rank}
\end{equation}
This lemma treats the possible nonfundamentality of $A_0$ in the underlying
Pell group. It is used with the verified strict hypothesis
\eqref{eq:p13-rank-hypotheses}; it is not replaced by informal cancellation
of $\Delta$ in the strong equation.

The auxiliary equation in \eqref{eq:p13-restored-roots}, the congruence
$V\equiv-c\pmod f$, and the odd abscissa quotient identity give an odd
auxiliary index $j_0$ with $j_0\equiv\pm p\pmod m$.
On the other hand $Q\equiv0\pmod c$, $f^2\equiv1\pmod c$, and
$V\equiv-R\pmod c$ give $j_0\equiv\pm R\pmod c$.
The signed index theorem from the ordinary source, Section~5, applies
because $0<2p<m$. Since $c\mid m$ and $0<p,R<c/2$, it follows that
$p=R$. This step permits either sign of $V$.

Now $n<p=R$, so $|2n-R-1|<R<\mu$, and the same congruence forces
$2n=R+1$. Consequently $r=k-2n$ is divisible by $4XY^2$, hence by both
$E$ and $U$. Set
\[
 h=\frac rE,\qquad \gamma_{\mathrm{old}}=\gamma-\frac rU.
\]
The first is a positive integer because $\psi_{P_0}(n)>n$ for $n\ge25$.
The second equals $(g-b_0)/H$. Since $n\ge2$, $k\ge4P_0>X$ and $c>X$;
furthermore $g>(a+1)c$, so $g-b_0>c-X>0$.
Thus every original ordinary compiler coordinate has been restored with
its required positive sign. The inherited compiler soundness gives the
blank K215 halt and hence a finite ZFC refutation.

\subsection{Constructive completeness}

Conversely, start with a canonical positive ordinary compiler witness of
a blank K215 halt. Canonical normalized witnesses supply these by
$i_{\mathrm{ordinary}}=\Delta i_{\mathrm{normalized}}$.
Its index theorem gives $p=R=2n-1$, so $r=k-2n$ is a positive multiple
of $4XY^2$ and therefore of $U$. Define
\[
 \gamma=\gamma_{\mathrm{old}}+\frac rU,
 \qquad t=\tau-\ell+y.
\]
Both are positive integers. In particular
$\tau^2=1+\ell(\ell+k)>\ell^2$ implies $\tau>\ell$.
Retain all other supplied coordinates. Then $x=\tau+y$,
$A_1=\tau^2$, and $E_{\mathrm{aux}}=dy^2$, so
\eqref{eq:p13-norm-identity} gives $N=0$.
Also $D_*=Ug$, giving $N_{\mathrm{main}}=0$, while $N_t=N_s=1$
are unchanged. Therefore \eqref{eq:p13-main} holds.
This proves both existential directions, with no additional cofinality
condition on an auxiliary index.

\subsection{Exact degree and finite verification}

All witnesses have degree one and the four ports degree zero. Put
\[
 q_0=\beta J,\qquad E_0=wsq_0^4,\qquad
 c_0=(\eta+\zeta)sq_0^3,\qquad \ell_0=E_0c_0.
\]
Their degrees are $1,6,5,11$, respectively. In the inner norm,
\[
 \deg d=14,\quad \deg V=7,\quad\deg E_{\mathrm{aux}}=28,
 \quad\deg A_1=22,\quad\deg\delta=12.
\]
The three terms of the final expression in
\eqref{eq:p13-norm-identity} have degrees $52,64,56$.
Only the middle one reaches degree $64$, with homogeneous part
\begin{equation}
 N_{64}=-4E_0^6c_0^4T^2f^6.
 \label{eq:p13-inner-head}
\end{equation}
In the efficient evaluator the squared core has degree at most $56$;
its correction $-4dA_1E_{\mathrm{aux}}$ has exact degree $64$.

The main residual has the exact identity
\begin{align*}
 N_{\mathrm{main}}
 &=U^2\bigl[2ac(X+H\gamma)+(X+H\gamma)^2-Hc^2-1\bigr]\\
 &\quad-2UHr(X+ac+H\gamma)+H^2r^2.
\end{align*}
Its three displayed terms have degrees $26$, at most $25$, and at most
$20$, so
\[
 (N_{\mathrm{main}})_{26}=128q_0^8E_0^2c_0\gamma.
\]
The other factors have leaders
\[
 (N_s)_{22}=i^2c_0^4,\qquad
 (N_t)_2=w(q_0-F-\alpha)-zq_0.
\]
Thus $\deg P_{13}\le2+26+22+64=114$.
On the full diagonal, where every witness is replaced by one
indeterminate, the four leading coefficients are
\[
 -2,\qquad 256\beta^{19},\qquad16\beta^{12},\qquad-64\beta^{36}.
\]
Their product is $2^{19}\beta^{67}\ne0$ on the fixed coefficient slice.
This proves exact degree $114$ without treating the ports as additional
variables or using a numerical experiment as an upper-degree proof.

The primary circuit uses $83$ binary gates, all live. The independent
standard-library checker imports no builder. It authenticates the whole
tuple and all source pins; independently verifies $41$ named formulas
against the literal circuit; proves the norm, centering, and main
cancellation identities; checks the universal weighted degree bound;
and evaluates six complete polynomial rays over three finite fields
using the actual fixed tuple. Normal and optimized Python runs agree.
These finite checks verify transcription, fixed data, and degree. They do
not replace the unbounded mathematical compiler and calibration proofs.

The original source50 absorber, ordinary strong-rank proof, compact native
adapter, and source59 comparison are preserved in the included upstream
directory. Their exact source paths, commits, Git blob identities and
SHA256 digests are recorded in \texttt{SOURCE\_PINS.json}. The complete
formal-assistant dependency tree of the inherited K215 endpoint is not
claimed to have been rebuilt in this work.

\section{A quartic equation in forty-two unrestricted integers}
\label{sec:quartic42}

The new equation in this section has \emph{forty-two unrestricted integer
unknowns}, exact total degree four, and twenty-seven quadratic residuals.
It improves the supplied forty-five-integer quartic without increasing its
degree. The four fixed coefficient ports are the complete, unchanged tuple
of that predecessor. In particular, the high-mask port is already shifted.
The construction uses the direct Pell/compiler route; no Jones or Sun
variable-reduction theorem is needed for this result.

\paragraph{Theorem.}
With the fixed coefficients specified in the accompanying binding, the
polynomial $Q_{42}$ defined below satisfies
\[
 \mathrm{ZFC}\ \vdash\
 \bigl[\neg\exists\boldsymbol{x}\in\mathbb Z^{42}
                \;Q_{42}(\boldsymbol{x})=0\bigr]
 \ \longleftrightarrow\ \operatorname{Con}(\mathrm{ZFC}).
\]
The argument is an ordinary arithmetic proof, composed with the expressly
retained native compiler and machine-calibration theorems. The new literal
polynomial has not received an end-to-end Lean or Rocq certificate. An
exported derivation in the repository's object-language ZFC calculus is
also distinct from the mathematical proof given here.

\subsection{The idea: prove the sign after recovering the rank}

Two positive coordinates in the predecessor, $\eta$ and $\zeta$, are each
represented by a ternary positive form. Those two representations use six
extra integer unknowns. We replace them by one representation of
$\eta\zeta>0$, which uses three. This initially proves only that the two
coordinates have the same sign.

An existing square scale supplies the missing information. Replacing
$w=u^2$ by $w=4u^2$ makes a modulus $e$ more than four times the decoded
index $R$. The Pell equations first recover a rank $p=R$ without deciding
the common sign. A negative common sign would then require a positive
integer smaller than $e$ to be divisible by $e$. Thus both coordinates
are positive after all. Canonical accepting witnesses have $w$ equal to
a sufficiently large dyadic square, so the factor four does not obstruct
completeness.

This ordering is essential. Invoking the predecessor's whole soundness
theorem before recovering the missing positive guards would be circular.
The local rank argument below is what permits those guards to be removed.

\subsection{Every variable and every residual}

Write the four fixed positive integer ports as
$\beta,K_0,M_C,M_F$, where $\beta=2^{1064921}-1$ and $M_F$ denotes the
\emph{already shifted} source mask. The complete coefficient binding, not
that exponent alone, specifies the instance. Let the thirty-three
arithmetic unknowns be
\begin{gather*}
 J,F,\alpha,z,f,h,r,v,w,\tau,\eta,\zeta,y,\gamma,o,j,\\
 N,E,P,c,a,X,D,C_2,L,F_2,V,g,I,B_*,\Delta,u,\rho.
\end{gather*}
Add nine guard unknowns $A_1,A_2,A_3,B_1,B_2,B_3,T_1,T_2,T_3$.
All forty-two range over $\mathbb Z$. In particular, $C_2,F_2,B_*$ and
$\Delta$ are independent supplied coordinates; their defining equations
are among the residuals. Define only the following abbreviations:
\[
 q=\beta J+1,\qquad C=q-F-\alpha,\qquad k=\eta+\zeta,
 \qquad \operatorname{Pos}(x,y,t)=1+x^2+y^2+t(t+1).
\]
These abbreviations add no quantified unknowns. The residuals are
\begin{align}
 R_1&=L-q^2+1, &
 R_2&=N-q(L+1), \nonumber\\
 R_3&=X-wq, &
 R_4&=a-E-(2v+1)N, \nonumber\\
 R_5&=E(X+1)-Xa, &
 R_6&=c-\eta-(a-E)k, \nonumber\\
 R_7&=P-k-E(c-\eta), &
 R_8&=\tau^2-P(P-k)-1, \nonumber\\
 R_9&=D-X-ac-\gamma(4a+3). &&
 \label{eq:q42-block1}
\end{align}
The next block is
\begin{align}
 R_{10}&=\Delta-a^2-4a-3, \nonumber\\
 R_{11}&=C_2-c^2, \nonumber\\
 R_{12}&=D^2-\Delta C_2-1, \nonumber\\
 R_{13}&=k-1-hE
       -\bigl[(q-F)(N-q)-CL+(M_C+qM_F)J\bigr], \nonumber\\
 R_{14}&=(K_0+w)C-F-z(q-1), \nonumber\\
 R_{15}&=F_2-f^2, \nonumber\\
 R_{16}&=V+c-fo, \nonumber\\
 R_{17}&=V-jc+k-1-hE, \nonumber\\
 R_{18}&=g+V^2-y^2.
 \label{eq:q42-block2}
\end{align}
Finally, put
\begin{align}
 R_{19}&=I-(2r+1)C_2, &
 R_{20}&=B_*-I^2, \nonumber\\
 R_{21}&=B_*-\Delta(F_2-1), &
 R_{22}&=y^2-B_*g-1, \nonumber\\
 R_{23}&=w-4u^2, &
 R_{24}&=q-2\rho^2, \nonumber\\
 R_{25}&=F\alpha-\operatorname{Pos}(A_1,A_2,A_3), &&\nonumber\\
 R_{26}&=FC-\operatorname{Pos}(B_1,B_2,B_3), &&\nonumber\\
 R_{27}&=\eta\zeta-\operatorname{Pos}(T_1,T_2,T_3). &&
 \label{eq:q42-block3}
\end{align}
The required single equation is
\begin{equation}
 \boxed{\displaystyle Q_{42}=\sum_{i=1}^{27}R_i^2=0.}
 \label{eq:q42-polynomial}
\end{equation}
Every residual has degree at most two in the supplied unknowns. The
coordinate $\tau$ occurs only in $R_8$, with term $\tau^2$, so the
coefficient of $\tau^4$ in $Q_{42}$ is exactly one. This proves exact
degree four on the authentic fixed-coefficient slice, independently of
any generic degree-propagation estimate.

The literal JSON uses \texttt{v\_scale}, \texttt{Baux} and
\texttt{q\_half\_root} for $v,B_*,\rho$, and named triples for the three
positive forms. The verifier expands all twenty-seven formulas above
and checks exact equality with that JSON. In particular, no division,
variable exponent or unspecified coefficient appears in the definition.

\subsection{Preliminary guards and restored arithmetic}

For every integer triple, $\operatorname{Pos}(x,y,t)>0$, since
$t(t+1)\ge0$. Conversely every positive integer $n$ has such a
representation. Indeed $4n-3\equiv1\pmod4$ is not of the forbidden form
$4^a(8b+7)$, so the three-square theorem represents it by three squares.
Exactly one square is odd. Writing that square as $(2t+1)^2$ and the
others as $(2x)^2,(2y)^2$ gives the claimed representation.

At an integer zero of \eqref{eq:q42-polynomial}, every residual vanishes.
The dyadic guard gives
\[
 q=\beta J+1=2\rho^2.
\]
Since $\beta\ge3$, $J<0$ would make $q<0$, while $J=0$ would require
$1=2\rho^2$. Thus $J>0$ and, for the bound source tuple, $q\ge16$.
The other two retained guards give $F\alpha>0$ and $FC>0$. Hence
$F,\alpha,C$ have the same nonzero sign; their sum is $q>0$, so
\begin{equation}
 F>0,\qquad \alpha>0,\qquad C>0.
 \label{eq:q42-simplex}
\end{equation}

Put $s=2v+1$, a nonzero odd integer, and set $Y=sq^3$ as a
\emph{proof abbreviation}. The first residuals give
\begin{gather*}
 N=q^3,\quad X=wq,\quad a=E+Y,\quad E=XY,\quad
 c=kY+\eta,\quad P=XY^2k+k,\\
 H=4a+3,\quad \Delta=(a+2)^2-1,\quad
 D=X+ac+\gamma H,\quad I=(2r+1)c^2.
\end{gather*}
For example $R_5=0$ and $a=E+Y$ imply
$E(X+1)=X(E+Y)$, hence $E=XY$ with no division or nonzero assumption.
Define the proof-only packing integer
\begin{equation}
 R=\bigl[(q-1)(q-F)+\alpha\bigr](q^2-1)
       +(M_C+qM_F)J.
 \label{eq:q42-packing}
\end{equation}
The retained equations then assert
\begin{gather}
 \tau^2-XY^2(XY^2+1)k^2=1,
 \qquad D^2-\Delta c^2=1, \label{eq:q42-pell1}\\
 I^2=\Delta(f^2-1),
 \qquad I^2V^2-(I^2-1)y^2=1, \label{eq:q42-pell2}\\
 k=R+1+hE,
 \qquad V=fo-c=jc-R. \label{eq:q42-index}
\end{gather}
The last two identities expose the quotient congruences directly; an
unproved quotient-integrality assertion is unnecessary.

The source packing bounds require only \eqref{eq:q42-simplex} and the
authentic mask bounds at this stage. With
$G=q\alpha+(q-1)C$, one has
\[
 q\le G\le q^2-q,
 \qquad 0<(M_C+qM_F)J<(q-1)(1+2q).
\]
Consequently
\begin{equation}
 q(q^2-1)<R<q^4-2,
 \qquad R-2>3q+1.
 \label{eq:q42-rbounds}
\end{equation}
No positivity assumption on $z,\eta,\zeta$ or $Y$ has been used in
obtaining these bounds.

\subsection{Bootstrap with an unknown common sign}

The new guard says $\eta\zeta>0$. Let
\[
 \sigma=\operatorname{sign}(\eta)=\operatorname{sign}(\zeta),
 \qquad \kappa=|k|>0,
 \qquad \theta=\frac{\eta}{k}\in(0,1).
\]
The fraction $\theta$ is used only in the proof. First $u\ne0$.
Otherwise $w=X=E=0$, $a=Y$ and $D=ac+\gamma H$. Reducing the main
Pell norm modulo $H$ would give $0\equiv1\pmod H$, since
$\Delta\equiv a^2\pmod H$. But $|Y|\ge q^3$ makes
$|H|=|4Y+3|>1$. This is impossible. Therefore
\begin{equation}
 w=4u^2\ge4,\quad X\ge4q>0,\quad |Y|\ge q^3,
 \quad e:=|E|\ge4q^4>4R+8.
 \label{eq:q42-modulus}
\end{equation}

Set $A=|a+2|$, $P_0=2XY^2+1$, and $c_0=|c|$. Then
\[
 A>e>R+2,\qquad P_0>A\ge2,\qquad \Delta=A^2-1.
\]
For negative $Y$, the first inequality follows from
$A=(X+1)|Y|-2>e$; for positive $Y$ it is immediate. Moreover
\[
 \frac{c_0}{\kappa}=
 \begin{cases}
 |Y|+\theta,&Y>0,\\
 |Y|-\theta,&Y<0,
 \end{cases}
\]
so $c_0>\kappa>0$ in either case.

Let $\chi_A(n),\psi_A(n)$ be the integer Pell sequences defined by
\[
 \chi_A(n)+\psi_A(n)\sqrt{A^2-1}
       =(A+\sqrt{A^2-1})^n.
\]
Ordinary Pell classification in \eqref{eq:q42-pell1} gives positive
ranks $n,p$ with
\begin{equation}
 \kappa=2\psi_{P_0}(n),\quad |\tau|=\chi_{P_0}(n),
 \qquad c_0=\psi_A(p),\quad |D|=\chi_A(p).
 \label{eq:q42-ranks}
\end{equation}
The first fundamental unit is the standard
$2b+1+2\sqrt{b(b+1)}$, where $b=XY^2>1$.
Since $P_0>A$ and $c_0>\kappa$, strict Pell growth gives $p>n$.
Also $P_0\equiv1\pmod e$, so the recurrence gives
\begin{equation}
 \sigma\,2n\equiv R+1\pmod e.
 \label{eq:q42-signedindex}
\end{equation}
If $\sigma=1$, its least positive representative gives $2n\ge R+1$.
If $\sigma=-1$, it gives $2n\ge e-R-1>R+1$. Thus
\[
 n\ge(R+1)/2\ge24,\qquad p>n.
\]
In particular the elementary Pell bounds imply
\begin{equation}
 c_0>A\Delta^2,\qquad c_0>2p,\qquad c_0>2R.
 \label{eq:q42-cbounds}
\end{equation}
For the first, $p\ge6$ already suffices, since
$\psi_A(6)=32A^5-32A^3+6A>A(A^2-1)^2$. The second follows by the
recurrence for $A\ge2,p\ge3$; the third follows from
$c_0\ge\psi_A(2)=2A>2R$.

\subsection{Recovering the rank without a sign assumption on \texorpdfstring{$c$}{c}}

We use the coefficient-independent ordinary strong-rank lemma retained
with the source proof. In its needed form, if
\[
 c_0=\psi_A(p)>A(A^2-1)^2,
 \qquad (i c_0^2)^2=(A^2-1)(f^2-1),\quad i>0,
\]
then there is a positive integer $m$ such that
\begin{equation}
 |f|=\chi_A(m),\quad p\mid m,\quad c_0\mid m,
 \quad m\ge c_0>2p,
 \quad i c_0^2=(A^2-1)\psi_A(m).
 \label{eq:q42-strong}
\end{equation}
Here its hypotheses hold with $i=|2r+1|$, using
\eqref{eq:q42-pell2} and \eqref{eq:q42-cbounds}.

For clarity, the two substantive ingredients in that inherited lemma are
strong divisibility and an elementary growth comparison. Write the
quadratic field's fundamental integer-coordinate unit using a parameter
$F_0$, and write $A=\chi_{F_0}(d)$,
$S=\psi_{F_0}(d)$, and
$\Delta=(F_0^2-1)S^2$. Classification yields
$|f|=\chi_{F_0}(t)$ and
$|I|=(F_0^2-1)S\psi_{F_0}(t)$.
Strong divisibility, together with $c_0^2\mid I$, gives
$\psi_{F_0}(\gcd(t,dp))\ge c_0/\Delta$.
If that gcd were a proper divisor of $dp$, doubling its index would give
\[
 2\psi_{F_0}(\gcd(t,dp))^2
   <\psi_{F_0}(2\gcd(t,dp))
   \le S c_0<A c_0,
\]
contrary to $c_0>A\Delta^2$. Hence $dp\mid t$, giving
$t=dm$ and $p\mid m$. Writing $m=pz$, the recurrence modulo $c_0^2$
gives $c_0\mid\Delta z$. Finally
$\gcd(c_0,\Delta)=\gcd(p,\Delta)$ follows from
$\psi_A(p)\equiv pA^{p-1}\pmod\Delta$, and gives $c_0\mid pz=m$.
The archived ordinary-kernel proof supplies the complete signed and
fundamental-unit details of this lemma.

The auxiliary norm in \eqref{eq:q42-pell2} excludes $V=0$ and $y=0$.
Pell classification at the integer parameter $|I|>1$ gives an odd rank
$\ell=2v_0+1$ such that
\[
 |I|\,|V|=\chi_{|I|}(\ell),\qquad
 |y|=\psi_{|I|}(\ell),\qquad |V|=Q_{v_0}(I^2).
\]
Oddness follows from the even-index residue
$\chi_{|I|}(2t)\equiv(-1)^t\pmod{|I|}$.
The standard odd-subsequence polynomials satisfy
\[
 Q_{v_0}(1-A^2)=(-1)^{v_0}\psi_A(\ell),
 \qquad Q_{v_0}(0)=(-1)^{v_0}\ell.
\]
They follow directly from $Q_0(x)=1$, $Q_1(x)=4x-3$ and
$Q_{t+1}(x)=(4x-2)Q_t(x)-Q_{t-1}(x)$.

The exposed identities in \eqref{eq:q42-index} imply
\[
 V\equiv-c\pmod{|f|},\qquad V\equiv-R\pmod{c_0}.
\]
As $I^2\equiv1-A^2\pmod{|f|}$, the first congruence gives
$\psi_A(\ell)\equiv\pm c_0\pmod{|f|}$, with all signs of $c,V$
absorbed in the ambiguity. Squaring and applying the Pell doubling formula,
\[
 \chi_A(2\ell)\equiv\chi_A(2p)\pmod{\chi_A(m)}.
\]
The signed step-down argument gives $\ell\equiv\pm p\pmod m$.
Indeed reduce $2\ell$ modulo $2m$, using
$\chi_A(t+2m)\equiv-\chi_A(t)$, then reflect indices above $m$.
The resulting representative is $\pm\chi_A(s)$, $0\le s\le m$.
The case $s=m$ gives residue zero. A negative representative is impossible
because
$\chi_A(s)+\chi_A(2p)\le2\chi_A(m-1)<\chi_A(m)$.
Strict growth then forces $s=2p$, proving the claimed index congruence.
Its hypotheses $0<2p<m$ were proved in \eqref{eq:q42-strong}.

Since $c_0\mid m$, this gives $\ell\equiv\pm p\pmod{c_0}$.
On the other hand $c_0\mid I$, so the second exposed congruence and
$Q_{v_0}(0)=(-1)^{v_0}\ell$ give
$\ell\equiv\pm R\pmod{c_0}$. Therefore
$p\equiv\pm R\pmod{c_0}$. The strict bounds
$0<p,R<c_0/2$ in \eqref{eq:q42-cbounds} force
\begin{equation}
 p=R.
 \label{eq:q42-mainrank}
\end{equation}

\subsection{The missing signs and the reverse construction}

If the common sign were negative, \eqref{eq:q42-signedindex} would make
$e$ divide $2n+R+1$. But $n<p=R$ and \eqref{eq:q42-modulus} give
\[
 0<2n+R+1<3R+1<e,
\]
a contradiction. Thus $\eta,\zeta>0$. Represent them separately by
$\operatorname{Pos}$, and set the old square coordinate to
$u_{\mathrm{old}}=2u$. Every residual of the authenticated predecessor
now vanishes. Its remaining signed-scale and compiler arguments therefore
produce a finite ZFC refutation. The missing individual guards have been
restored before invoking that whole-system theorem.

There is a small transcription correction in the predecessor's written
transport discussion. Its literal circuit, retained in $R_{14}$, has
\[
 (K_0+w)C=F+z(q-1),
\]
whereas the prose writes $z-1$. The needed conclusion is $z>0$.
After its Pell recovery, $w>q>F$ and $C\ge1$ make the left side greater
than $F$, hence $z(q-1)>0$. This proves the required guard for the
literal shifted coordinate. Neither the equation nor its coefficients
are changed by correcting that sentence.

Conversely, a finite refutation supplies the predecessor's canonical
accepting witnesses. Its complete coefficient adapter ensures
\[
 q=2^t,\qquad t\ \hbox{odd},\qquad R\ \hbox{odd},
 \qquad w=2^{R-t},\qquad R>3q.
\]
Thus $R-t$ is even and at least four. Choose
\[
 u=2^{(R-t-2)/2},
\]
so $w=4u^2$. The canonical $\eta,\zeta$ are positive; represent their
product by one $\operatorname{Pos}$ triple and keep the other coordinates.
All twenty-seven new residuals vanish. This is an existence construction
using canonical witnesses, rather than an invertible substitution on
every predecessor tuple.

\subsection{What has been checked}

The new construction retains twenty-five predecessor residuals exactly,
strengthens the square scale, and replaces two positive guards by one.
It changes $(45,28,4)$ to $(42,27,4)$ in the order
unknowns, residuals, and single-polynomial degree.
With the four ports left symbolic, the exact expanded sum of squares
has 391 terms, compared with 407 for the predecessor. This is a symbolic
description count, not a claim about coefficient bit complexity or an
optimal arithmetic circuit.

The main builder and a separate implementation expand the literal
residuals in sparse integer polynomial arithmetic. They check variable
counts and occurrence, the unchanged rows, the two modifications, exact
degree four, and the coefficient one of $\tau^4$. Both work under
ordinary and optimized Python. A separate transcription check verifies
all twenty-seven factored formulas above against the emitted JSON.
The parent polynomial and its proof are authenticated by their pinned
Git blob hashes; the complete coefficient tuple and semantic source
dependencies are retained with their own provenance.

These finite checks establish the displayed algebra and source binding.
The universal solvability equivalence rests on the two-way mathematical
argument and its expressly identified inherited arithmetic and compiler
theorems. No consistency assumption, refutation search, global
minimality assertion or external peer review is being substituted for
that proof.

\section{A sextic in thirty-five integers and an octic in thirty}
\label{sec:companions}

The same fixed coefficient tuple also gives two equations with fewer
unrestricted integer unknowns:
\[
 \begin{array}{c|r|r|r}
  \text{equation}&\text{unknowns}&\text{residuals}&\text{exact degree}\\\hline
  S_{35}&35&20&6\\
  O_{30}&30&15&8
 \end{array}
\]
Both satisfy the same ZFC-provable equivalence between absence of integer
zeros and $\operatorname{Con}(\mathrm{ZFC})$. The collection's historical
sextic and octic entries use forty-two and thirty-six unrestricted
integers, respectively. These are improvements of seven and six
unknowns relative to those entries, without a global minimality claim.

\subsection{Move a square guard to the coordinate that will be eliminated}

Let $Q_{42}^{X}$ have the same residuals as $Q_{42}$ except that
\[
 q-2\rho^2=0\qquad\text{is replaced by}\qquad X-8\xi^2=0.
\]
This replaces one unknown by one unknown. The retained equations include
$w=4u^2$ and $X=wq$.

\paragraph{Square-guard lemma.}
The integer zero sets of $Q_{42}$ and $Q_{42}^{X}$ are in bijection under
$\xi=u\rho$, with all other coordinates fixed.

\paragraph{Proof.}
The equations $F\alpha>0$ and $FC>0$ imply that $F,\alpha,C$ have the
same nonzero sign. Thus $q=F+\alpha+C$ satisfies $|q|\ge3$, before either
square guard has been used. If $u=0$, then $w=X=0$, the row
$E(X+1)=Xa$ gives $E=0$, and
\[
 a=(2v+1)q^3,\qquad |a|\ge27.
\]
With $H=4a+3$, the main norm has $D=ac+\gamma H$ and
$\Delta=a^2+4a+3\equiv a^2\pmod H$. Hence
$D^2-\Delta c^2=1$ implies $H\mid1$, contrary to $|H|\ge105$.
Therefore $u\ne0$.

At a zero of $Q_{42}^{X}$, the scale equations give
$qu^2=2\xi^2$. Write $\xi/u=A/B$ in lowest terms with $B>0$.
Then $qB^2=2A^2$, so $B^2\mid2$ and $B=1$. Consequently
$\rho=\xi/u$ is an integer and $q=2\rho^2$, restoring the deleted
guard. Conversely, a zero of $Q_{42}$ has
$X=4u^2\cdot2\rho^2=8(u\rho)^2$. The two substitutions are inverse.
\hfill$\square$

This elementary relocation matters for degree: $X=8\xi^2$ is quadratic,
while direct substitution of $w=4u^2$ into $X=wq$ would give a cubic
expression for $X$. The primary quartic remains the one in
Section~\ref{sec:quartic42}.

\subsection{The thirty-five-integer sextic, in full}

Keep the abbreviations $q=\beta J+1$, $C=q-F-\alpha$,
$k=\eta+\zeta$, and put $s=2v+1$. The twenty-six arithmetic unknowns are
\begin{gather*}
 J,F,\alpha,z,f,h,r,v,\tau_*,\eta,\zeta,y,\gamma,o,j,\\
 E,c,a,D,L,V,I,B_*,\Delta,u,\xi.
\end{gather*}
Together with the nine guard unknowns from Section~\ref{sec:quartic42},
there are thirty-five unknowns, all ranging over $\mathbb Z$.
Here $D,L,B_*,\Delta$ remain independent coordinates. The twenty
residuals are
\begin{align}
 S_1&=L-q^2+1, &
 S_2&=8\xi^2-4u^2q, \nonumber\\
 S_3&=a-E-sq(L+1), &
 S_4&=E(8\xi^2+1)-8\xi^2a, \nonumber\\
 S_5&=c-\eta-(a-E)k, &&\nonumber\\
 S_6&=\tau_*^2+(2\tau_*-k)E(c-\eta)-1, &&\nonumber\\
 S_7&=D-8\xi^2-ac-\gamma(4a+3), &&\nonumber\\
 S_8&=\Delta-a^2-4a-3, &
 S_9&=D^2-\Delta c^2-1.
 \label{eq:sextic35-first}
\end{align}
The other residuals are
\begin{align}
 S_{10}&=k-1-hE
    -\bigl[((q-1)(q-F)+\alpha)L+(M_C+qM_F)J\bigr], \nonumber\\
 S_{11}&=(K_0+4u^2)C-F-z(q-1), \nonumber\\
 S_{12}&=V+c-fo, &
 S_{13}&=V-jc+k-1-hE, \nonumber\\
 S_{14}&=I-(2r+1)c^2, &
 S_{15}&=B_*-I^2, \nonumber\\
 S_{16}&=B_*-\Delta(f^2-1), &
 S_{17}&=y^2-B_*(y^2-V^2)-1, \nonumber\\
 S_{18}&=F\alpha-\operatorname{Pos}(A_1,A_2,A_3), &&\nonumber\\
 S_{19}&=FC-\operatorname{Pos}(B_1,B_2,B_3), &&\nonumber\\
 S_{20}&=\eta\zeta-\operatorname{Pos}(T_1,T_2,T_3). &&
 \label{eq:sextic35-second}
\end{align}
The equation is
\begin{equation}
 \boxed{\displaystyle S_{35}=\sum_{i=1}^{20}S_i^2=0.}
 \label{eq:sextic35}
\end{equation}
Every residual has degree at most three. Only $S_{14}$ contains $r$,
and the coefficient of $r^2c^4$ in $S_{35}$ is exactly four. Thus its
degree is exactly six after the authentic coefficients are substituted.

To derive this system, first change the root coordinate by
\[
 \tau_{\mathrm{old}}=\tau_*+P-k,
\]
an invertible integral linear change. Then eliminate the seven registers
\[
 \begin{array}{c|l}
 P&k+E(c-\eta)\\
 F_2&f^2\\
 g&y^2-V^2\\
 X&8\xi^2\\
 N&q(L+1)\\
 C_2&c^2\\
 w&4u^2.
 \end{array}
\]
Each register has a defining equation with coefficient $1$ or $-1$.
After substitution, the old root restores as
$\tau_{\mathrm{old}}=\tau_*+E(c-\eta)$. The first norm becomes
\[
 [\tau_*+E(c-\eta)]^2
 -[k+E(c-\eta)]E(c-\eta)-1
 =\tau_*^2+(2\tau_*-k)E(c-\eta)-1.
\]
The cancellation of the fourth-degree term is exact. Every zero of the
twenty-residual system uniquely restores a zero of $Q_{42}^{X}$, and
every zero of $Q_{42}^{X}$ yields one in the other direction.

\subsection{The thirty-integer octic, in full}

The twenty-one arithmetic unknowns are now
\[
 J,F,\alpha,z,f,h,r,v,\tau,\eta,y,\gamma,o,j,E,c,a,V,I,u,\xi.
\]
Add the same nine guard unknowns. All thirty range over $\mathbb Z$.
The following are polynomial abbreviations, not independent unknowns:
\begin{gather*}
 q=\beta J+1,\quad C=q-F-\alpha,\quad s=2v+1,\quad H=4a+3,\\
 k=1-V+jc+hE,\quad b=E(c-\eta),\quad \Delta=a^2+4a+3,\\
 R=\bigl[(q-1)(q-F)+\alpha\bigr](q^2-1)+(M_C+qM_F)J.
\end{gather*}
The fifteen residuals are
\begin{align}
 O_1&=8\xi^2-4u^2q, &
 O_2&=a-E-sq^3, \nonumber\\
 O_3&=E(8\xi^2+1)-8\xi^2a, &
 O_4&=c-\eta-(a-E)k, \nonumber\\
 O_5&=\tau^2-(k+b)b-1, &&\nonumber\\
 O_6&=(8\xi^2+ac+\gamma H)^2-\Delta c^2-1, &&\nonumber\\
 O_7&=-V+jc-R, &
 O_8&=(K_0+4u^2)C-F-z(q-1), \nonumber\\
 O_9&=V+c-fo, &
 O_{10}&=I-(2r+1)c^2, \nonumber\\
 O_{11}&=I^2-\Delta(f^2-1), &
 O_{12}&=y^2-I^2(y^2-V^2)-1, \nonumber\\
 O_{13}&=F\alpha-\operatorname{Pos}(A_1,A_2,A_3), &&\nonumber\\
 O_{14}&=FC-\operatorname{Pos}(B_1,B_2,B_3), &&\nonumber\\
 O_{15}&=\eta(k-\eta)-\operatorname{Pos}(T_1,T_2,T_3). &&
 \label{eq:octic30-residuals}
\end{align}
The single equation is
\begin{equation}
 \boxed{\displaystyle O_{30}=\sum_{i=1}^{15}O_i^2=0.}
 \label{eq:octic30}
\end{equation}
Every residual has degree at most four. In $O_{11}$ the term
$-a^2f^2$ has coefficient $-1$. No other residual contains both $a$ and
$f$, so the coefficient of $a^4f^4$ in $O_{30}$ is exactly one. The
degree is therefore exactly eight on the fixed-coefficient slice.

Here the old root coordinate $\tau$ is retained. The twelve eliminated
coordinates restore, after composing the substitutions, as follows:
\[
 \begin{array}{c|l@{\qquad}c|l}
 N&q^3 & C_2&c^2\\
 F_2&f^2 & g&y^2-V^2\\
 w&4u^2 & X&8\xi^2\\
 \zeta&k-\eta & B_*&I^2\\
 D&8\xi^2+ac+\gamma H & \Delta&a^2+4a+3\\
 P&k+b & L&q^2-1.
 \end{array}
\]
The actual monic elimination order is
$N,C_2,F_2,g,w,X,\zeta,B_*,D,\Delta,P,L$.
In particular, $\zeta$ is first removed using
$\zeta=P-E(c-\eta)-\eta$, and $P$ is later removed using the exposed
congruence row. This yields the quadratic abbreviation for $k$, so the
common-sign guard still has degree only three.

Each elimination is an exact integral bijection on zero sets. The square
guard lemma then restores the original $Q_{42}$ system. Applying its
theorem proves the claimed consistency equivalences for both companions.

\paragraph{Reproducibility and scope.}
The companion generator checks nineteen monic eliminations, all fifty-four
parent-row restoration identities, all thirty-five factored residual
transcriptions, and both exact-degree coefficient certificates. It uses
explicit exceptions rather than removable Python assertions. Normal and
optimized runs agree. The literal and expanded JSON packets preserve the
whole coefficient binding; no new coefficient transfer is assumed.
These finite algebra checks supplement the ordinary proof. They do not
constitute a proof-assistant certificate for the inherited unbounded
Pell and compiler theorems. The JSON uses \texttt{rho} for $\xi$ and
\texttt{tau\_tail} for $\tau_*$.

\section{A compact nine-positive companion}
\label{sec:compact-nine}

The supplied companion polynomial has nine strictly positive unknowns,
exact total degree $1\,834\,516$, and a live arithmetic circuit of $453$
binary additions, subtractions and multiplications.  At the inspected
ProveIt revision, package59's corresponding point is
$(9,1\,910\,404,473)$.  Thus this component saves $75\,888$ degrees and
$20$ operations.  It also binds the source to the complete compact44
coefficient tuple, whose four unsigned binary encodings occupy
$295\,388$ bytes, compared with $533\,099$ for the preceding tuple.
These are distinct measures: the constants remain large, and operation
count is not bit complexity.  This is a local Pareto improvement, not a
claim of optimality or a major advance in universal variable counts.

The authoritative equation is the ordinary integer polynomial represented
by the supplied arithmetic DAG:
\begin{center}
\small\texttt{nine/generated/split\_116\_positive\_global.json}.
\end{center}
Its fixed
leaves are completely materialized and authenticated by
\texttt{nine/fixed/BINDING.json}.  Intermediate register names below are
abbreviations, not quantified variables.  In particular, all displayed
exponents are fixed integers.  The nine inputs, in order, are
\[
 b,g,h,s,w,C,i,j,n\in\mathbb Z_{>0}.
\]
Using the identified inherited native compiler, its compact whole-tuple
adapter, the calibration from K215 to finite refutations and the Jones ratio theorem,
ZFC proves
\[
 (\exists b,g,h,s,w,C,i,j,n>0\;\mathcal F=0)
 \quad\Longleftrightarrow\quad \neg\operatorname{Con}(\mathrm{ZFC}).
\]
The added reduction is ordinary mathematics formalizable in ZFC.  The
package supplies exact finite certificates, not a new end-to-end
Lean/Rocq proof or an exported object-language ZFC derivation.

\subsection{The source specialization and exact target certificate}

Use the $31$-positive-coordinate, $22$-row source in
\texttt{nine/data/parent\_system60.json}.  Its four symbolic coefficient
ports are $\beta,K_0,M_C,M_F$.  The complete compact44 tuple has bit
lengths
\[
 (590069,592884,590069,590070),
\]
with aliases $K_0=K$ and $M_F=M_{F,\mathrm{source}}$; the last port is
already shifted by the native base minus one.  The four payloads are
supplied in \texttt{nine/coefficients/}.  The supplied native builder
reconstructs them exactly from the literal transition data.  The finite
verification includes all $66062$ native basis identities and the
explicit rank controller.

The old package59 source must be adjusted before this tuple is used.
With all four ports symbolic, exact expansion gives
\[
 r_{15}^{\mathrm{new}}=r_{15}^{\mathrm{old}}-\beta J,
 \qquad r_i^{\mathrm{new}}=r_i^{\mathrm{old}}\quad(i\ne15).
\]
This is the justified transport shift $z_{\mathrm{old}}=z+1$; its
completeness uses the inherited source bound $z_{\mathrm{old}}\ge2$ on
zeros.  It is not an arbitrary positive-coordinate substitution.  Apply
the invertible row operations
\[
 \rho_{14}=r_{14}+r_2,\qquad \rho_{15}=r_{15}+r_3,
 \qquad \rho_i=r_i\ (i\notin\{14,15\}).
\]
The resulting source has $89$ records: $7$ constant, $39$ linear and
$43$ quadratic.  In its original variable order choose weights
\[
\begin{split}
(w_i)={}&(1,7,61,36,16,85,20,67,43,11,34,13,96,80,29,50,\\
 &2,109,110,112,116,91,53,55,27,101,92,9,107,37,60).
\end{split}
\]
Assign target exponents in the following paired order:
\[
\begin{split}
\text{rows}:\;&(2,15,8,17,13,18,12,20,21,4,22,9,11,19,1,6,10,16,3,14,7,5),\\
\text{targets}:\;&(28,124,325,423,598,658,775,904,941,955,1027,1236,\\
 &1245,1313,1323,1530,1571,1584,1637,1787,2014,2036).
\end{split}
\]
Write $q_j$ for the target of row $j$, and set
$U_y(T)=\sum_i y_iT^{w_i}$.  For a coefficient $a_{j,\alpha}$ of
$y^\alpha$ in $\rho_j$, let $d=|\alpha|$ and
$e(\alpha)=\sum_i\alpha_iw_i$.  Define
\[
 c_{j,\alpha}=
 \begin{cases}
 2a_{j,\alpha},&d=0,1,\\
 a_{j,\alpha}\prod_i\alpha_i!,&d=2,
 \end{cases}
 \qquad
 D_d(T)=\sum_{j,\,|\alpha|=d}
       c_{j,\alpha}T^{q_j-e(\alpha)}.
\]
The supplied literal sparse kernels and target certificate prove
\[
 [T^{q_j}]\bigl(D_2(T)U_y(T)^2+D_1(T)U_y(T)+D_0(T)\bigr)
       =2\rho_j(y) \quad (1\le j\le22).
\]
Indeed, among all $22544$ same-degree source/expansion interactions,
only the $89$ intended contributions hit selected targets.  Squares
have multiplicity one in $U_y^2$ and cross terms multiplicity two, so
the stated weighting recovers twice the source coefficient.  This is a
formal coefficient identity, with arbitrary witnesses and arbitrary
coefficient ports.  The resulting degrees are
\[
 (\deg D_0,\deg D_1,\deg D_2)=(1571,1934,1836),\quad W=116,
 \quad H=\max_d(\deg D_d+dW)=2068.
\]
The leading coefficient of $D_2$ is exactly $+1$.  There are no negative
kernel exponents.  The maximal target is $2036$, leaving the required
headroom below the convolution height.

\subsection{Source words and the exact binary threshold}

Replacing native-port powers by their strict bit bounds and summing
absolute values gives the sharp common coefficient majorant
$M=2^{1182954}$: every actual weighted coefficient has magnitude strictly
less than $M$.  Set
\[
 \kappa=1182966,\quad a=2^\kappa=4096M,\quad B=ab^2,
 \quad \ell=\sum_iB^{w_i},\quad U_s=\ell+g.
\]
Define the following integer polynomials:
\[
\begin{aligned}
 G_s&=(D_2(B)U_s+D_1(B))U_s+D_0(B),&
 T_*&=\sum_jB^{q_j},\\
 N_0&=B^{117},&Z&=B^{2069},\\
 \mathrm{Mask}&=N_0-1-(b-1)\ell,&
 J_*&=2G_s+T_*,\\
 S_*&=g+N_0J_*,&
 T_{**}&=\mathrm{Mask}+N_0(B-2)T_*.
\end{aligned}
\]
For every $b,g\ge1$, we have $U_s\ge B^{116}$ and
$D_2(B)>B^{1836}/2$.  The two potentially negative contributions, relative
to $B^{1836}U_s^2$, are bounded by
\[
 \frac{|D_1U_s|}{B^{1836}U_s^2}<\frac{2M}{B^{18}}<\frac18,
 \qquad
 \frac{|D_0|}{B^{1836}U_s^2}<\frac{2M}{B^{497}}<\frac18.
\]
Thus $G_s>0$ unconditionally; also $0\le\mathrm{Mask}<N_0$, so the source
words are nonnegative before any finalizer consequence is used.
On the full soundness range $g<4bB^{116}$, one has $U_s<6bB^{116}$ and
\[
 \frac{2G_s}{Z}<\frac{86Mb^2}{B-1}<\frac1{32},\qquad
 T_*<\frac Z{64},\quad(B-2)T_*<\frac Z{64},\quad g<N_0.
\]
Consequently $S_*+T_{**}<B^{2186}$.  For binary $b$, the low mask forces
$g=\sum_i z_iB^{w_i}$ with $0\le z_i<b$, so $U_s$ encodes the positive
witnesses $y_i=z_i+1$.  Every convolution coefficient is bounded in
magnitude by $993Mb^2<B/4$.  The inherited signed-borrow padding argument,
now with the certified target identities and these numerical bounds,
therefore gives
\[
 S_*\mathbin{\&}T_{**}=0
 \quad\Longleftrightarrow\quad
 g\text{ encodes a positive common source zero}.
\]
A positive source zero supplies such a word for every sufficiently large
binary $b$, with $0<g<bB^{116}$.  The strict positivity follows because
the all-one source assignment has $r_5=1$.

Now put
\[
 L=bB^{2186},\quad
 R=L(S_*+T_{**}+1)-S_*-1,\quad
 N=2^{8872234}B^{1155}.
\]
Unconditionally $R,N\ge8$ and $b\le R,N$.  If $b=2^k$, write
$v=1182966+2k$ and $u=k+2186v$, so $B=2^v$ and $L=2^u$.
The bounded source words give the two-block expansion
\[
 R=(S_*+T_{**})L+(L-1-S_*).
\]
With $\operatorname{pop}$ denoting binary population count and
$\tau=\operatorname{pop}(S_*)+\operatorname{pop}(T_{**})
-\operatorname{pop}(S_*+T_{**})$, this implies
\[
 \operatorname{pop}(R)=u+\operatorname{pop}(T_{**})-\tau.
\]
The low mask contributes $117v-31k$ bits and the $22$ high digits
contribute $22v-22$ bits.  Hence the exact calibration is
\[
 u+\operatorname{pop}(T_{**})=2325v-30k-22
 =2(8872234+1155v)=2\log_2N.
\]
Since $\tau\ge0$, with equality exactly for binary-disjoint words,
Legendre's identity
$v_2\binom{2R}{R}=\operatorname{pop}(R)$ proves
\[
 N^2\mid\binom{2R}{R}
 \quad\Longleftrightarrow\quad S_*\mathbin{\&}T_{**}=0.
\]
The case $k=0$ is included and cannot produce a final zero, since then
$\mathrm{Mask}=N_0-1$ and $0<g<N_0$.

\subsection{The explicit final equation}

The remaining formulas are the inherited free-coordinate finalizer,
applied with the parameters just proved.  Write $j_0=j-1\ge0$ and define
\[
\begin{aligned}
 U&=N^2w,&Y&=N^2s,&M_0&=RY,\\
 P&=2M_0^2U,&K&=R+1+h(P-1),&A&=M_0(U+1),\\
 B_1&=2R+1,&W_1&=bw,&J&=4A-5,\\
 Q&=3W_1C-2(W_1^2-1),&&&D&=(A^2-1)C^2+1,\\
 k_0&=iJC,&E&=k_0DC,&F_1&=(A^2-1)E^2+1,\\
 G&=A+F_1(CD+1-A),&&&H_1&=C+F_1(B_1+j_0C),\\
 I&=(G^2-1)H_1^2+1,&a_0&=(P^2-1)K^2+1,&a_1&=DF_1I.
\end{aligned}
\]
Finally set
\[
\begin{aligned}
 t&=bB^{116},&m&=K^2(4t-g)-16t(C-KY)^2,\\
 V_0&=PK,&V_1&=k_0(8D^2-4D-1)GH_1,\\
 \Omega&=Jm,&E_0&=J(n+W_1^2)+Q,\\
 \Theta&=E_0-\Omega(V_0+V_1).
\end{aligned}
\]
The equation is
\begin{equation}
 \boxed{\mathcal F=
 [\Theta^2-\Omega^2(a_0+64a_1)]^2
       -256\Omega^4a_0a_1=0.}
 \label{eq:compact-nine}
\end{equation}
The inherited proof applies because $R,N\ge8$, $b\le R,N$, the source
words are unconditionally nonnegative, and the source completeness bound
$g<bB^{116}$ has the required factor-$32$ margin.
The important positivity step occurs before square-root integrality:
\[
 V_0^2-a_0=K^2-1>0,
 \quad
 V_1^2-64a_1=
 [(8D+1)k_0^2-64D](GH_1)^2+64DF_1(H_1^2-1)>0.
\]
Also $E_0>0$.  A norm zero with $m=0$ would give $E_0^4>0$;
otherwise its real factorization forces $m>0$.  This yields
$|C-KY|<K/2$ and the needed carrier bound.  The rational-root separation
then makes both square roots integral, gives $J\mid Q$, and invokes the
Jones ratio theorem to recover binary $b$, the source range and binomial
divisibility.  Conversely the source and ratio witnesses yield an integer
$n>0$ by the explicit bound $W_1^2+Q/J<2D<V_1$.  Thus no unproved cofinality
claim about $j$ or $n$ enters the construction.  The complete inherited
argument and the new applicability checks are retained in
\texttt{nine/PROOF.md} and its pinned upstream dependencies.

\subsection{Exact degree, cost, and verification}

The source heads give
\[
 \deg S_*=4370,\quad\deg T_{**}=4308,
 \quad r=\deg R=8743,\quad q=\deg N=2310.
\]
Let $\Delta_1=V_1^2-64a_1$.  To resolve the apparent leading cancellation
in the norm, use the exact identity
\[
 \Theta^2-\Omega^2(a_0+64a_1)
 =E_0^2-2E_0\Omega(V_0+V_1)
  +\Omega^2\bigl[(K^2-1)+\Delta_1+2V_0V_1\bigr].
\]
The separate degree certificate has a positive diagonal leading term and
no unresolved positive-degree cancellation.  The diagonal reaches the
multivariate upper bound, proving the exact degree
\[
 104r+400q+8W+316=1\,834\,516.
\]
The initial Horner circuit has $467$ operations.  Its $102$ pure powers of
$B$ are replaced by a shared $88$-gate addition chain.  Every chain row
$e=a+b$ certifies $B^e=B^aB^b$, and exact structural normalization checks
the complete rewritten circuit, including the definition of $B$.  The
final tally is $162$ additions, $22$ subtractions and $269$
multiplications, totaling $453$.  Scalar multiplication and every power
multiplication are included; constant leaves and reuse are free.

Replacing every positive input by a nonnegative input plus one gives the
supplied nine-nonnegative companion, with the same degree and $460$
operations.  No nine-variable unrestricted-integer claim is made.
The $12$ final nonliteral leaves occupy $2\,217\,161$ raw magnitude bytes
or $47\,364$ deterministic gzip bytes; the signs and hashes are explicit.

The command \texttt{python3 nine/reproduce.py} reconstructs the entire
component offline using the Python standard library.  It regenerates the
source and both domains, verifies all target identities, reconstructs the
native tuple, materializes the derived leaves, and checks the power-chain
and exact-degree certificates.  A second finite checker imports no
project builder or verifier.  All authored checks use exceptions and
remain enabled under \texttt{python3 -O}.  Modular and small binary-block
fixtures are separately labeled as finite transcription checks; they do
not establish the unbounded consistency semantics.


\section{Verification, source correction, and reproducibility}
\label{sec:verification}

The archive contains the final polynomial definitions, their fixed
coefficient bindings, the mathematical proofs, and offline verification
programs. Its new checks use the Python standard library. Historical
scripts that require other packages are preserved only as source evidence
when relevant; the new replay does not need to execute them. The root
command is
\begin{verbatim}
python3 verify_all.py
\end{verbatim}
Run it from the extracted package directory. It authenticates the package,
executes the component checks, reconstructs the bound native coefficient
tuples, and replays the finite polynomial and circuit certificates. The
receipt is \path{verification/REPLAY.json}. Both ordinary and optimized
Python execution were used during preparation; acceptance does not rely
on removable \texttt{assert} statements.

For $Q_{42}$, two separate sparse polynomial implementations check the
literal residuals, the $45\to42$ modification, all variable occurrences,
the exact degree, and the coefficient one of $\tau^4$. A transcription
check compares the displayed factored equations with the JSON. The
coefficient replay regenerates all four source-34 numerals and checks
$69\,230$ native basis conditions and $8\,801$ coefficient digits. The
companions have explicit restoration maps after the proved scale-guard
exchange: deleting coordinates through
monic equations preserves the integer solution set by substitution, with
no rational divisions.

For $P_{13}$, the certificate separates a multivariate degree upper bound
from a nonzero homogeneous leading value, so a sampled univariate degree
is not mistaken for an upper bound. The norm identity, unit-product
residues, complete operation ledger, and fixed coefficient binding are
checked. The nine-variable replay verifies all $22\,544$ sparse target
interactions, the coherent native source change, exact degree data,
every shared power identity, and the complete circuit rewrite. It checks
all $66\,062$ native basis identities of the compact tuple. These are
finite certificates for exact algebra and data. They do not replace the
unbounded two-way compiler and Pell arguments.

One source correction is necessary. The predecessor quartic's prose
writes the transport coordinate as $z-1$, whereas its authenticated
literal equation is
\[
 (K_0+w)(q-F-\alpha)=F+z(q-1).
\]
The correct required restored guard is $z>0$. The old positive product
coordinate would be $z+1$. After the source Pell recovery,
$w>q>F$ and $q-F-\alpha\ge1$ prove this strict positivity directly.
The new quartic proof states the correction; the original source file is
retained unchanged for comparison. Neither the literal polynomial nor
its coefficients are silently repaired to match the prose.

The reviews included in the archive are internal independent-agent
mathematical or implementation reviews in this research session. They
are identified as such. They are not external human peer review, and
finite diagnostics are labeled separately from universal proofs.

\subsection{File map}

\begin{longtable}{@{}p{0.31\textwidth}p{0.65\textwidth}@{}}
\toprule
Path & Purpose\\
\midrule
\endhead
\path{report.tex}, \path{report.pdf} & This article and its editable source.\\
\path{positive13/} & The degree-$114$ polynomial, complete norm proof,
fixed tuple binding, degree and circuit certificates.\\
\path{quartic42/} & All 27 quadratic residuals of the quartic,
common-sign proof, exact expansion and independent algebra check.\\
\path{quartic42/companions/} & The sextic and octic, coordinate
restoration data and exact elimination certificates.\\
\path{quartic42/dependencies/} & Authenticated ordinary compiler and
Pell dependencies, source-34 coefficient payloads and their replay.\\
\path{nine/} & The optimized nine-variable circuits, sparse code and
kernel certificates, compact native tuple and final fixed leaves.\\
\path{RESULTS.json} & Machine-readable domain, variable, degree and
operation ledger for the new forms.\\
\path{SOURCE_SNAPSHOT.json} & Comparison revision and source locations.\\
\path{MANIFEST.json} & SHA256 hashes of delivered files.\\
\path{verify_all.py} & One-command offline finite verification.\\
\path{RESEARCH_NOTES.md} & Unproved directions and rejected shortcuts,
kept separate from the theorems.\\
\bottomrule
\end{longtable}

Compile the article with
\begin{verbatim}
pdflatex -interaction=nonstopmode -halt-on-error report.tex
pdflatex -interaction=nonstopmode -halt-on-error report.tex
\end{verbatim}
The PDF was rendered and inspected after compilation. The source contains
the complete formulas rather than relying on remote images or an online
notebook. The archive includes the upstream licenses and notices; no
external repository or publication is modified by this contribution.

\section{Where a larger advance could come from}
\label{sec:research}

The new reductions suggest specific further work, but they do not yet
resolve the request for a major change of scale. The two clearest new
mechanisms are reusable: recover a missing sign from an already available
rank modulus, and center a root before eliminating it by a norm.
Both exploit constraints already paid for elsewhere in the equation.

\paragraph{Recover more signs without duplicating positive forms.}
The quartic argument saves three witnesses because its common-sign
ambiguity is a single bit that the index congruence can decide. A further
reduction would need a comparable argument for other guards. Removing
them on the strength of the old whole-system theorem would be circular:
that theorem assumes those very guards. In particular, simply allowing
both scale coordinates to be negative changes the main binomial recovery
and has not been justified here.

\paragraph{Use a smaller direct binary-disjointness interface.}
The nine-variable companion spends most of its degree converting
\emph{binary disjointness} of two packed nonnegative integers $S,T$ into
a divisibility threshold on a central binomial coefficient. The direct
condition
\[
 \binom{S+T}{S}\equiv1\pmod2
\]
expresses absence of binary carries. A sufficiently economical
Diophantine interface for this condition could avoid the large two-block
index and threshold used here. The carry interpretation by itself is
not such an interface; a new explicit polynomial construction with a
complete converse is still needed.

\paragraph{Combine root elimination with divisibility, not only radicals.}
Norm equations are attractive because they remove witnesses, but a
rational square root need not satisfy a required quotient congruence.
The thirteen-variable norm works because the first recovered root is
monic over the integers and the other root is then an integer difference.
Eliminating additional roots while losing one of the divisibility guards
can admit extraneous solutions. For example,
\[
 63\left((4/3)^2-1\right)=7^2
\]
shows that the strong Pell-looking equation alone does not force a
rational $f$ to be integral. This is a counterexample to that proposed
local shortcut, not to the full guarded source system. No smaller
equation obtained by that unsupported cancellation is claimed.

\paragraph{Change the proof-checking source itself.}
A different primitive-recursive refutation checker, direct arithmetic
verification of a proof certificate, or a different finite tiling language
could reduce coefficient and compiler complexity. An appeal to MRDP alone
does not supply the necessary literal coefficients, degree bounds, or
two-way compiler proof. This package does not replace K215 by an
unimplemented alternative.

The contributions ready for repository review are the explicit equations
and reductions proved above. Further priority, global optimality, and
proof-assistant certification require separate work.

\begin{thebibliography}{99}
\raggedright
\bibitem{matiyasevich}
Yu.~V. Matiyasevich,
\emph{The Diophantineness of enumerable sets},
Dokl. Akad. Nauk SSSR \textbf{191}(2) (1970), 279--282;
English translation, Soviet Math. Doklady \textbf{11} (1970), 354--357.
\url{https://www.mathnet.ru/eng/dan35274}.

\bibitem{jones}
J.~P. Jones, \emph{Universal Diophantine equation},
Journal of Symbolic Logic \textbf{47}(3) (1982), 549--571.
\url{https://doi.org/10.2307/2273588}.

\bibitem{sun2026}
Z.-W. Sun, \emph{Ten unknowns for Hilbert's tenth problem over the integers},
arXiv:2609.01594v2, 2 September 2026, preprint.
\url{https://arxiv.org/abs/2609.01594v2}.

\bibitem{repo57}
ProveIt contributors, \emph{Separated degrees}, package 57,
including \path{quartic/QUARTIC45_PROOF.md} and the low-degree source.
Pinned comparison commit \snapshot.
\url{https://github.com/VladimirReshetnikov/ProveIt/tree/beb0192cb8dc36f78c369c080f2adfbc1887a486/Computability/DiophantineConsistency/con-zfc-diophantine-forms/57-separated-degrees}.

\bibitem{repo58}
ProveIt contributors, \emph{Integral root}, package 58,
pinned comparison commit \snapshot.
\url{https://github.com/VladimirReshetnikov/ProveIt/tree/beb0192cb8dc36f78c369c080f2adfbc1887a486/Computability/DiophantineConsistency/con-zfc-diophantine-forms/58-integral-root}.

\bibitem{repo59}
ProveIt contributors, \emph{Degree separated}, package 59,
including the ordinary merger, thirteen-positive adjunct, and nine-free
adjunct, pinned comparison commit \snapshot.
\url{https://github.com/VladimirReshetnikov/ProveIt/tree/beb0192cb8dc36f78c369c080f2adfbc1887a486/Computability/DiophantineConsistency/con-zfc-diophantine-forms/59-degree-separated}.

\bibitem{repodeps}
ProveIt contributors, inherited coefficient, ordinary Pell, signed-scale,
native compiler and K215 calibration sources. Exact historical revisions,
repository paths, byte hashes and local reading map are supplied in
\path{quartic42/dependencies/SOURCE_PINS.json},
\path{quartic42/dependencies/README.md}, and the component
\path{SOURCE_PINS.json} files. Original licenses and notices accompany
the authenticated copies.
\end{thebibliography}
\end{document}
